% Источники: exam/materials/преподаватель/2020/14-04-2020-Model-L9-1-14.04.20.pdf; exam/materials/преподаватель/2020/14-04-2020-Model-L9-2-14.04.20.pdf; exam/materials/моделирование.pdf, раздел 19.
\section{Аналитические модели систем массового обслуживания.}

\textbf{Система массового обслуживания (СМО)} — система, функционирование которой состоит в обработке поступающих заявок. Аналитическая модель СМО описывает вероятностное поведение системы уравнениями теории массового обслуживания и позволяет получить характеристики в стационарном режиме.

\subsection*{Нотация Кендалла $A/B/c$}

\begin{itemize}
  \item $A$ — закон распределения интервалов между поступлениями заявок;
  \item $B$ — закон распределения времени обслуживания;
  \item $c$ — число каналов обслуживания.
\end{itemize}

Символы: \textbf{M} (Markovian) — пуассоновский поток / экспоненциальный закон; \textbf{D} — детерминированное время; \textbf{G} — общий закон распределения. В расширенной записи указывают ёмкость системы, размер источника и дисциплину обслуживания.

\subsection*{Коэффициент загрузки}

\[
  a = \lambda b = \frac{\lambda}{\mu}, \qquad \rho = \frac{a}{c} = \frac{\lambda}{c\mu},
\]
где $\lambda$ — интенсивность входного потока, $b$ — среднее время обслуживания, $\mu=1/b$ — интенсивность обслуживания одним каналом, $c$ — число каналов. Условие устойчивости для системы с ожиданием:
\[
  \lambda < c\mu \quad \text{или} \quad \rho < 1.
\]

\subsection*{Простейшая СМО M/M/1}

Для одноканальной системы с пуассоновским потоком ($\lambda$) и экспоненциальным обслуживанием ($\mu = 1/b$) условие стационарности $\rho = \lambda/\mu < 1$. Из уравнений Колмогорова в стационарном режиме:

\[
  p_0 = 1-\rho, \qquad p_k = (1-\rho)\rho^k,\quad k=0,1,\ldots
\]
\[
  L_s = \frac{\rho}{1-\rho}, \qquad
  W_s = \frac{1}{\mu - \lambda},
\]
\[
  L_q = \frac{\rho^2}{1-\rho}, \qquad
  W_q = \frac{\rho}{\mu - \lambda},
\]
где $L_s$, $L_q$ — среднее число заявок в системе и в очереди; $W_s$, $W_q$ — среднее время пребывания в системе и в очереди.

Связь средних характеристик задаётся формулами Литтла:
\[
  L_s = \lambda W_s, \qquad L_q = \lambda W_q.
\]

\subsection*{Виды СМО}

\begin{itemize}
  \item \textbf{Открытая} — источник заявок расположен вне системы.
  \item \textbf{Закрытая} — заявка после обслуживания возвращается в поток.
  \item \textbf{СМО без ожидания (отказная)} — заявка, заставшая все каналы занятыми, покидает систему.
  \item \textbf{СМО с ограниченной очередью} — очередь не превышает заданного размера.
  \item \textbf{СМО с неограниченной очередью} — любая заявка ждёт освобождения канала.
\end{itemize}

\subsection*{СМО с отказами M/M/c/c}

Если очередь отсутствует и заявка при занятости всех $c$ каналов теряется, стационарные вероятности имеют вид:
\[
  p_k = \frac{a^k/k!}{\sum_{i=0}^{c} a^i/i!}, \qquad k=0,\ldots,c.
\]
Вероятность отказа равна вероятности занятости всех каналов:
\[
  P_{\text{отк}} = p_c =
  \frac{a^c/c!}{\sum_{i=0}^{c} a^i/i!}.
\]
Относительная пропускная способность $Q = 1-P_{\text{отк}}$, абсолютная пропускная способность $A = \lambda Q$.

\subsection*{Аналитическая модель vs.\ имитационная}

Аналитическая модель СМО строится на основе теории вероятностей и позволяет получить характеристики системы в \textbf{стационарном режиме} в замкнутом виде (при определённых допущениях о распределениях). Имитационная модель снимает ограничения на распределения и позволяет исследовать переходные режимы, однако требует значительно больших вычислительных ресурсов.

\clearpage
